\documentclass[a4paper,12pt]{article}
\usepackage[utf8]{inputenc}
\usepackage[T1]{fontenc}
\usepackage[ngerman]{babel}
\usepackage{amsmath,amssymb}
\usepackage{booktabs}
\usepackage{enumitem}
\usepackage[margin=2.5cm]{geometry}
\usepackage{pdflscape}
\usepackage{tcolorbox}

\tcbuselibrary{breakable}

\newtcolorbox{merkbox}[1][]{
  colback=blue!5,
  colframe=blue!40,
  fonttitle=\bfseries,
  title={Merke},
  #1
}

\newtcolorbox{analogie}[1][]{
  colback=green!5,
  colframe=green!40,
  fonttitle=\bfseries,
  title={Analogie},
  #1
}

\begin{document}

\begin{center}
{\LARGE\bfseries \"Ubungsblatt 1 -- Erkl\"arungen}\\[0.5em]
{\large VC Logik}
\end{center}

\bigskip

%% ============================================================
%% AUFGABE 1
%% ============================================================
\section{Aufgabe 1: Vollst\"andig geklammerte Formeln}

\subsection{Was wird hier verlangt?}

Wir bekommen Formeln \emph{ohne} (oder mit wenigen) Klammern und sollen
sie so umschreiben, dass \textbf{jeder} bin\"are Operator sein eigenes
Klammerpaar bekommt. Dadurch wird die Struktur der Formel eindeutig --
man braucht keine Vorrangregeln mehr.

\begin{analogie}
Denk an Mathematik: $2 + 3 \times 4$ ist nicht eindeutig, wenn man die
Regel "`Punkt vor Strich"' nicht kennt. Die vollst\"andig geklammerte
Version $\bigl(2 + (3 \times 4)\bigr)$ macht sofort klar, was zuerst
berechnet wird.

In der Logik funktioniert es genauso -- nur mit anderen Operatoren.
\end{analogie}

\subsection{Die Schl\"usselidee: Operatorpr\"azedenz}

Genau wie "`Punkt vor Strich"' gibt es in der Aussagenlogik eine feste
Reihenfolge, welcher Operator am st\"arksten bindet:

\begin{center}
\renewcommand{\arraystretch}{1.3}
\begin{tabular}{clp{7cm}}
\toprule
\textbf{Rang} & \textbf{Operator} & \textbf{Bedeutung} \\
\midrule
1 (st\"arkste) & $\neg$ (Negation)       & Bindet immer zuerst -- wie ein Minus vor einer Zahl \\
2              & $\wedge$ (Und)          & Wie Multiplikation \\
3              & $\vee$ (Oder)           & Wie Addition \\
4              & $\to$ (Implikation)     & "`Wenn \ldots, dann \ldots"' \\
5 (schw\"achste) & $\leftrightarrow$ (\"Aquivalenz) & "`Genau dann, wenn \ldots"' \\
\bottomrule
\end{tabular}
\end{center}

\begin{merkbox}
\textbf{Faustregel:} Klammere immer von innen nach au\ss en --
starte beim st\"arksten Operator ($\neg$) und arbeite dich zum
schw\"achsten ($\leftrightarrow$) vor.
\end{merkbox}

\subsection{L\"osungen Schritt f\"ur Schritt}

\subsubsection*{(a) $A \vee \neg\neg B \wedge \neg\neg\neg C$}

\begin{enumerate}[leftmargin=2cm]
\item[\textbf{Schritt 1:}] \textbf{Negationen erkennen} ($\neg$ bindet am st\"arksten).

$\neg\neg B$ bedeutet: Erst $\neg B$, dann nochmal $\neg$ davor.\\
$\neg\neg\neg C$ bedeutet: $\neg C$, dann $\neg\neg C$, dann $\neg\neg\neg C$.

Die Negationen sind hier schon eindeutig -- sie brauchen keine
zus\"atzlichen Klammern, weil $\neg$ immer auf die \emph{n\"achste}
Formel wirkt.

\item[\textbf{Schritt 2:}] \textbf{$\wedge$ klammern} (bindet st\"arker als $\vee$).

$\neg\neg B \wedge \neg\neg\neg C$ wird zu
$(\neg\neg B \wedge \neg\neg\neg C)$.

\item[\textbf{Schritt 3:}] \textbf{$\vee$ klammern}.

$A \vee (\neg\neg B \wedge \neg\neg\neg C)$ wird zu
$(A \vee (\neg\neg B \wedge \neg\neg\neg C))$.
\end{enumerate}

\begin{merkbox}
Ergebnis: $\boxed{(A \vee (\neg\neg B \wedge \neg\neg\neg C))}$
\end{merkbox}

\subsubsection*{(b) $A \vee \neg C \leftrightarrow \neg B \wedge C$}

\begin{enumerate}[leftmargin=2cm]
\item[\textbf{Schritt 1:}] \textbf{Negationen:} $\neg C$ und $\neg B$ --
bereits eindeutig.

\item[\textbf{Schritt 2:}] \textbf{$\wedge$:}
$\neg B \wedge C$ wird zu $(\neg B \wedge C)$.

\item[\textbf{Schritt 3:}] \textbf{$\vee$:}
$A \vee \neg C$ wird zu $(A \vee \neg C)$.

\item[\textbf{Schritt 4:}] \textbf{$\leftrightarrow$:}
$(A \vee \neg C) \leftrightarrow (\neg B \wedge C)$ wird zu
$((A \vee \neg C) \leftrightarrow (\neg B \wedge C))$.
\end{enumerate}

\begin{merkbox}
Ergebnis: $\boxed{((A \vee \neg C) \leftrightarrow (\neg B \wedge C))}$
\end{merkbox}

\subsubsection*{(c) $A \leftrightarrow \neg C \to B \wedge C$}

\begin{enumerate}[leftmargin=2cm]
\item[\textbf{Schritt 1:}] \textbf{Negation:} $\neg C$ -- eindeutig.

\item[\textbf{Schritt 2:}] \textbf{$\wedge$:}
$B \wedge C$ wird zu $(B \wedge C)$.

\item[\textbf{Schritt 3:}] \textbf{$\to$:}
$\neg C \to (B \wedge C)$ wird zu $(\neg C \to (B \wedge C))$.

\item[\textbf{Schritt 4:}] \textbf{$\leftrightarrow$:}
$A \leftrightarrow (\neg C \to (B \wedge C))$ wird zu
$(A \leftrightarrow (\neg C \to (B \wedge C)))$.
\end{enumerate}

\begin{merkbox}
Ergebnis: $\boxed{(A \leftrightarrow (\neg C \to (B \wedge C)))}$
\end{merkbox}

\subsubsection*{(d) $\neg A \leftrightarrow B \to \neg C \wedge A \vee B$}

\begin{enumerate}[leftmargin=2cm]
\item[\textbf{Schritt 1:}] \textbf{Negationen:} $\neg A$ und $\neg C$ -- eindeutig.

\item[\textbf{Schritt 2:}] \textbf{$\wedge$:}
$\neg C \wedge A$ wird zu $(\neg C \wedge A)$.

\item[\textbf{Schritt 3:}] \textbf{$\vee$:}
$(\neg C \wedge A) \vee B$ wird zu $((\neg C \wedge A) \vee B)$.

\item[\textbf{Schritt 4:}] \textbf{$\to$:}
$B \to ((\neg C \wedge A) \vee B)$ wird zu
$(B \to ((\neg C \wedge A) \vee B))$.

\item[\textbf{Schritt 5:}] \textbf{$\leftrightarrow$:}
$\neg A \leftrightarrow (B \to ((\neg C \wedge A) \vee B))$ wird zu
$(\neg A \leftrightarrow (B \to ((\neg C \wedge A) \vee B)))$.
\end{enumerate}

\begin{merkbox}
Ergebnis: $\boxed{(\neg A \leftrightarrow (B \to ((\neg C \wedge A) \vee B)))}$
\end{merkbox}

\newpage

%% ============================================================
%% AUFGABE 2
%% ============================================================
\section{Aufgabe 2: Teilformeln}

\subsection{Was ist eine Teilformel?}

\begin{analogie}
Stell dir eine Formel wie ein \textbf{Lego-Bauwerk} vor. Die Teilformeln
sind alle Einzelteile und Zwischenprodukte, aus denen das Bauwerk
zusammengebaut wird -- vom kleinsten Einzelstein (den Variablen $A$, $B$, $C$)
bis hin zum fertigen Geb\"aude (der gesamten Formel).
\end{analogie}

Formaler: Eine Teilformel ist jede Formel, die als "`Baustein"' beim
schrittweisen Aufbau der Gesamtformel entsteht.

\subsection{Vorgehensweise}

Am einfachsten findet man alle Teilformeln, indem man die Formel
\textbf{von innen nach au\ss en aufbaut}:

\begin{enumerate}
\item Starte mit den \textbf{Atomen} (den Variablen).
\item Wende Schritt f\"ur Schritt Operatoren an.
\item Jedes Zwischenergebnis ist eine Teilformel.
\end{enumerate}

\subsection{Anwendung auf die Formel}

Formel: $\neg A \leftrightarrow B \to \neg\neg\neg C \wedge A \vee \neg\neg B$

\medskip
Vollst\"andig geklammert:
$(\neg A \leftrightarrow (B \to ((\neg\neg\neg C \wedge A) \vee \neg\neg B)))$

\bigskip
\textbf{Schritt-f\"ur-Schritt-Aufbau:}

\begin{center}
\renewcommand{\arraystretch}{1.4}
\begin{tabular}{clp{7.5cm}}
\toprule
\textbf{Nr.} & \textbf{Teilformel} & \textbf{Erkl\"arung} \\
\midrule
1  & $A$   & Atom (Variable) \\
2  & $B$   & Atom (Variable) \\
3  & $C$   & Atom (Variable) \\
\midrule
4  & $\neg C$         & Negation von $C$ \\
5  & $\neg\neg C$     & Negation von Teilformel 4 \\
6  & $\neg\neg\neg C$ & Negation von Teilformel 5 \\
\midrule
7  & $\neg B$         & Negation von $B$ \\
8  & $\neg\neg B$     & Negation von Teilformel 7 \\
\midrule
9  & $\neg A$         & Negation von $A$ \\
\midrule
10 & $(\neg\neg\neg C \wedge A)$
   & Konjunktion von Tf.\,6 und Tf.\,1 \\
11 & $((\neg\neg\neg C \wedge A) \vee \neg\neg B)$
   & Disjunktion von Tf.\,10 und Tf.\,8 \\
12 & $(B \to ((\neg\neg\neg C \wedge A) \vee \neg\neg B))$
   & Implikation von Tf.\,2 nach Tf.\,11 \\
13 & $(\neg A \leftrightarrow (B \to ((\neg\neg\neg C \wedge A) \vee \neg\neg B)))$
   & \"Aquivalenz von Tf.\,9 und Tf.\,12 \\
\bottomrule
\end{tabular}
\end{center}

\begin{merkbox}
Insgesamt gibt es \textbf{13~Teilformeln}. Die Gesamtformel selbst
(Nr.\,13) ist immer auch eine Teilformel.
\end{merkbox}

\newpage

%% ============================================================
%% AUFGABE 3
%% ============================================================
\section{Aufgabe 3: Wahrheitstabelle}

\subsection{Was ist eine Wahrheitstabelle?}

\begin{analogie}
Eine Wahrheitstabelle ist wie ein \textbf{Fahrplan}: F\"ur jede
m\"ogliche Kombination der Eingaben (wahr/falsch) zeigt sie dir
systematisch das Ergebnis. Nichts wird \"ubersehen, weil \emph{alle}
M\"oglichkeiten durchprobiert werden.
\end{analogie}

\subsection{Wie viele Zeilen braucht man?}

Wir haben 3~Variablen ($A$, $B$, $C$). Jede kann entweder 0~(falsch)
oder 1~(wahr) sein. Also gibt es $2^3 = 8$~m\"ogliche Kombinationen
= 8~Zeilen.

\subsection{Wie f\"ullt man die Tabelle aus?}

\begin{enumerate}
\item \textbf{Eingabespalten} ($A$, $B$, $C$): Alle $2^3 = 8$ Kombinationen
      systematisch auflisten (wie bin\"ar z\"ahlen: 000, 001, 010, \ldots).
\item \textbf{Einfache Spalten zuerst}: Die Negationen ($\neg C$,
      $\neg\neg C$, \ldots) kann man direkt aus den Eingaben ablesen.
\item \textbf{Zusammengesetzte Spalten}: Von links nach rechts ausf\"ullen --
      jede Spalte braucht nur die Werte der vorherigen Spalten.
\end{enumerate}

\begin{merkbox}
\textbf{Erinnerung an die Wahrheitswerte der Operatoren:}

\smallskip
\begin{tabular}{ll}
$\neg A$: & Kehrt den Wert um (aus 0 wird 1, aus 1 wird 0) \\
$A \wedge B$: & Nur wahr, wenn \textbf{beide} wahr sind \\
$A \vee B$: & Wahr, wenn \textbf{mindestens einer} wahr ist \\
$A \to B$: & Nur falsch, wenn $A$ wahr und $B$ falsch ist \\
$A \leftrightarrow B$: & Wahr, wenn beide den \textbf{gleichen} Wert haben \\
\end{tabular}
\end{merkbox}

\subsection{Beispielzeile im Detail}

Nehmen wir $A=1,\; B=0,\; C=0$:

\begin{enumerate}[leftmargin=2cm]
\item[\textbf{Neg.:}]
$\neg C = \neg 0 = 1$, \quad
$\neg\neg C = \neg 1 = 0$, \quad
$\neg\neg\neg C = \neg 0 = 1$

$\neg B = \neg 0 = 1$, \quad
$\neg\neg B = \neg 1 = 0$

$\neg A = \neg 1 = 0$

\item[$F_{10}$:]
$\neg\neg\neg C \wedge A = 1 \wedge 1 = 1$

\item[$F_{11}$:]
$F_{10} \vee \neg\neg B = 1 \vee 0 = 1$

\item[$F_{12}$:]
$B \to F_{11} = 0 \to 1 = 1$
\quad (Implikation mit falscher Pr\"amisse ist \emph{immer} wahr!)

\item[$F_{13}$:]
$\neg A \leftrightarrow F_{12} = 0 \leftrightarrow 1 = 0$
\quad (unterschiedliche Werte $\Rightarrow$ falsch)
\end{enumerate}

\subsection{Vollst\"andige Wahrheitstabelle}

Abk\"urzungen:
\begin{align*}
F_{10} &= (\neg\neg\neg C \wedge A) &
F_{11} &= (F_{10} \vee \neg\neg B) &
F_{12} &= (B \to F_{11}) &
F_{13} &= (\neg A \leftrightarrow F_{12})
\end{align*}

\begin{center}
\renewcommand{\arraystretch}{1.25}
\begin{tabular}{ccc|cccccc|cccc}
\toprule
$A$ & $B$ & $C$
  & $\neg C$ & $\neg\neg C$ & $\neg\neg\neg C$
  & $\neg B$ & $\neg\neg B$ & $\neg A$
  & $F_{10}$ & $F_{11}$ & $F_{12}$ & $F_{13}$ \\
\midrule
0 & 0 & 0  & 1 & 0 & 1  & 1 & 0 & 1  & 0 & 0 & 1 & \textbf{1} \\
0 & 0 & 1  & 0 & 1 & 0  & 1 & 0 & 1  & 0 & 0 & 1 & \textbf{1} \\
0 & 1 & 0  & 1 & 0 & 1  & 0 & 1 & 1  & 0 & 1 & 1 & \textbf{1} \\
0 & 1 & 1  & 0 & 1 & 0  & 0 & 1 & 1  & 0 & 1 & 1 & \textbf{1} \\
1 & 0 & 0  & 1 & 0 & 1  & 1 & 0 & 0  & 1 & 1 & 1 & \textbf{0} \\
1 & 0 & 1  & 0 & 1 & 0  & 1 & 0 & 0  & 0 & 0 & 1 & \textbf{0} \\
1 & 1 & 0  & 1 & 0 & 1  & 0 & 1 & 0  & 1 & 1 & 1 & \textbf{0} \\
1 & 1 & 1  & 0 & 1 & 0  & 0 & 1 & 0  & 0 & 1 & 1 & \textbf{0} \\
\bottomrule
\end{tabular}
\end{center}

\subsection{Was sagt uns die Tabelle?}

\begin{merkbox}
Schau dir die letzte Spalte ($F_{13}$) an: Sie ist \textbf{immer 1,
wenn $A = 0$} und \textbf{immer 0, wenn $A = 1$}. Das bedeutet:

Die gesamte Formel verh\"alt sich genauso wie $\neg A$ -- sie ist
logisch \"aquivalent zu $\neg A$.

\medskip
Die Formel ist weder eine Tautologie (nicht immer wahr) noch eine
Kontradiktion (nicht immer falsch) -- sie ist \textbf{erf\"ullbar}.
\end{merkbox}

\end{document}
